\chapter{Materiais e metodoloxía}
\label{chap:materiais-metodoloxia}

Neste capítulo descríbense os materiais empregados e a metodoloxía seguida para o desenvolvemento
do sistema proposto. En particular, preséntase o conxunto de datos utilizado, os pasos de
preprocesado aplicados ás imaxes, a estrutura xeral do pipeline deseñado e a configuración
experimental empregada para avaliar o comportamento do sistema.

\section{Conxunto de datos}
O \textit{dataset} empregado para este traballo foi cedido polo  \textcolor{red}{\textbf{TODO: Preguntar qué departamento en específico nolo cedeu.}}.
Este conxunto componse de 4704 imaxes de algas \gls{kelp}, que, como xa se mencionou en **METER REFERENCIA A ACHEGAMENTO DO PROBLEMA**, teñen unha 
amplia variabilidade de captura. 

\section{Etiquetado e variables de interese}
\textcolor{red}{\textbf{TODO: Explicar o significado de HPI e IVR, a interpretación dos seus valores
e a súa relevancia para o problema estudado.}}

\section{Preprocesado das imaxes}
\textcolor{red}{\textbf{TODO: Detallar os pasos de normalización, limpeza e transformación aplicados
ás imaxes antes da inferencia.}}

\section{Arquitectura do pipeline}
\textcolor{red}{\textbf{TODO: Presentar a estrutura xeral do sistema, indicando as etapas principais,
como a detección da alga, o tratamento da rexión de interese e a estimación das variables obxectivo.}}

\section{Modelos empregados}
\textcolor{red}{\textbf{TODO: Describir os modelos utilizados en cada etapa do pipeline, incluíndo a
motivación da súa elección.}}

\section{Configuración experimental}
\textcolor{red}{\textbf{TODO: Explicar a partición dos datos, as métricas empregadas e o procedemento
seguido para avaliar o sistema.}}
